\documentclass[11pt,a4paper]{article}
\usepackage[margin=27mm]{geometry}
\usepackage{longtable,booktabs,array,fvextra}
\usepackage{amsmath,amsthm,amssymb,fontspec,unicode-math,hyperref}
\setmainfont{STIXGeneral}
\setmathfont{STIX Two Math}
\setmonofont{Menlo}
\hypersetup{hidelinks}
\usepackage{newunicodechar}
\newfontfamily\numeralfallback{STIX Two Text}
\newunicodechar{∎}{\ensuremath{\blacksquare}}
\newunicodechar{⁰}{{\numeralfallback\char"2070\relax}}
\newunicodechar{₀}{{\numeralfallback\char"2080\relax}}
\newunicodechar{₁}{{\numeralfallback\char"2081\relax}}
\newunicodechar{₂}{{\numeralfallback\char"2082\relax}}
\newunicodechar{₃}{{\numeralfallback\char"2083\relax}}
\newunicodechar{₄}{{\numeralfallback\char"2084\relax}}
\newunicodechar{⁻}{{\numeralfallback\char"207B\relax}}
\newunicodechar{⁸}{{\numeralfallback\char"2078\relax}}
\newunicodechar{ⁱ}{{\numeralfallback\char"2071\relax}}
\newunicodechar{⁴}{{\numeralfallback\char"2074\relax}}
\newunicodechar{⁷}{{\numeralfallback\char"2077\relax}}
\newunicodechar{⁶}{{\numeralfallback\char"2076\relax}}
\newunicodechar{⁵}{{\numeralfallback\char"2075\relax}}
\newunicodechar{₈}{{\numeralfallback\char"2088\relax}}
\newunicodechar{ᵢ}{\textsubscript{i}}
\newunicodechar{ⱼ}{\textsubscript{j}}
\newunicodechar{ₜ}{\textsubscript{t}}
\providecommand{\tightlist}{\setlength{\itemsep}{0pt}\setlength{\parskip}{0pt}}
\newtheorem{theorem}{Theorem}
\newtheorem{proposition}[theorem]{Proposition}
\newtheorem{lemma}[theorem]{Lemma}
\newtheorem{corollary}[theorem]{Corollary}
\theoremstyle{remark}\newtheorem{remark}[theorem]{Remark}
\newcommand{\gen}[1]{\langle #1\rangle}
\newcommand{\normalgen}[1]{\langle\!\langle #1\rangle\!\rangle}
\newcommand{\Z}{\mathbb Z}
\newcommand{\F}{\mathbb F}
\setlength{\emergencystretch}{2em}

\providecommand{\passthrough}[1]{#1}

\title{The common quotient of finite two-group string markings}
\author{Human Compression Project}
\date{}
\begin{document}
\maketitle
\begin{abstract}
The forest quotient theorem identifies the third lower-central quotient of
every finite two-group string C-group whose consecutive markers do not
commute. It is the same binary path group in every case of a fixed rank.
This proves the common-quotient and regular-polytope covering theorems of
Hou and Schulte. We separately calculate the auxiliary Frattini-intersection
subgroup proposed in that source: even in the common quotient itself it is
generally nontrivial. This distinguishes the valid principal conclusions
from the source's incorrect prescribed-kernel formulas.
\end{abstract}

\section{The fixed path group}
A string C-group of rank $d$ has a generating list of nonidentity involutions
$r_0,\ldots,r_{d-1}$, satisfies $[r_i,r_j]=1$ for $|i-j|>1$, and has
the intersection property
\[
 G_J\cap G_K=G_{J\cap K},\qquad G_J=\gen{r_j:j\in J}.
\]
Nondegeneracy means that consecutive markers do not commute. For a finite
two-group, their product orders are then powers of two at least four.
Let $\gamma_3(G)=[G',G]$. All groups considered in the main theorem are
finite, with no bound on their nilpotency class.

For $d\ge2$, define $Q_d$ on $\F_2^d\times\F_2^{d-1}$ by
\[
 (v,z)(w,t)=(v+w,z+t+\beta(v,w)),\qquad
 \beta_i(v,w)=v_iw_{i-1}\quad(1\le i\le d-1).              \tag{1}
\]
The marker $r_j$ has vertex coordinate $e_j$ and zero edge coordinates.
The complete proof in UA.Forest shows that these markers generate $Q_d$,
that its order is $2^{2d-1}$, and that its marked subgroup on $J$ has
exactly the vertex coordinates in $J$ and the edge coordinates internal
to $J$. Thus its marking has IP and all Boolean deletions. Its consecutive
product squares are the distinct basis elements $c_i$ of the central edge
space, so its type is $\{4,\ldots,4\}$. At $d=2$ it is the dihedral
group of order eight.

This is also the group described as $G(C_d)$ in the source. One way to see
the identification without a group census is to use the common presentation:
involutory markers, distant commutations, adjacent products of order dividing
four, and central adjacent commutators. Collecting the markers followed by
the central commutators gives at most $2^{2d-1}$ elements; (1) realizes all
of them. The prescribed marking, not only the abstract group order, is fixed.
The regular-polytope correspondence gives the polytope $C_d$ attributed to
Conder in [HS26].

\section{The two principal source conclusions}
\begin{theorem}[Hou--Schulte, Theorems 1.1 and 1.2]
Every finite nondegenerate string C-group $G$ of two-power order and rank
$d\ge2$ has a generator-preserving epimorphism onto $Q_d$. In fact
\[
 G/\gamma_3(G)\cong Q_d                                  \tag{2}
\]
as marked groups. Every corresponding regular polytope covers $C_d$.
\end{theorem}
\begin{proof}
The noncommutation graph of the marking of $G$ is exactly the path on its
$d$ vertices: string commutation removes every other pair, and
nondegeneracy retains each path edge. In particular it is a forest. The
original marking has IP and all markers are nonidentity involutions, so
every hypothesis of UA.Forest holds. That theorem proves independence of
the edge commutators in $G'/\gamma_3(G)$, then identifies the entire
marked quotient with the vertex-and-edge model (1). It gives (2) and the
required marked epimorphism. The standard marked-cover correspondence
of [MS02], in the precise form stated in [UA], converts this epimorphism
into the asserted rank- and adjacency-preserving polytope cover.
\end{proof}

The quotient is independent of $G$ once the rank is fixed. It is itself a
finite nondegenerate string C-group, so it supplies a least element in the
marked covering order of the stated class. This is the source's principal
minimal-covering assertion. It is not a classification of all larger groups
or all fibers above this quotient.

The proof of UA.Forest is a complete part of the accompanying umbrella.
It retains Burnside's basis theorem, nilpotency of finite two-groups and
the standard lower-central graded Lie construction. Its separation argument
uses IP in the original group and proves independence before identifying
the quotient. No assumption that IP automatically survives quotienting,
no bound on nilpotency class and no finite catalogue enter the proof.

\section{The auxiliary kernel in the source}
The source's Theorem 4.2 gives a specific subgroup $C$ as the proposed
kernel. Its formulas do not hold in general. We keep the source's
definition and compute it in $Q_d$.
Put
\[
 \sigma_i=r_{i-1}r_i,\qquad
 A_i=\gen{\sigma_1,\ldots,\sigma_{i-1},\sigma_i^2,
                    \sigma_{i+1},\ldots,\sigma_{d-1}},
 \qquad C=\bigcap_{i=1}^{d-1}\Phi(A_i).                    \tag{3}
\]
These are equations (7)--(9) in the source. In the coordinate model,
$\sigma_i^2=c_i$. The vertex images of $\sigma_1,\ldots,\sigma_{d-1}$
are independent and span the even-weight vertex subspace. Every $A_i$
contains every central $c_j$: for $j\ne i$ it contains $\sigma_j^2$,
and $c_i$ is a displayed generator. Its vertex image has dimension $d-2$,
so $|A_i|=2^{2d-3}$.

For a finite two-group, Burnside's theorem gives
$\Phi(A_i)=A_i^2[A_i,A_i]$. Since our group is class two, the squares
and pair commutators of the displayed generators generate this subgroup.
At the two endpoints we obtain
\[
 \Phi(A_1)=\gen{c_2,\ldots,c_{d-1}},\qquad
 \Phi(A_{d-1})=\gen{c_1,\ldots,c_{d-2}}.                  \tag{4}
\]
All displayed coordinates are squares. The missing endpoint coordinate
cannot occur in another square or commutator: every remaining rotation
has zero vertex coordinate at that endpoint, while $c_i$ is central of
order two. This proves equality in (4), not just containment.
For an internal $i$, the other rotations supply every square $c_j$ with
$j\ne i$, and
$[\sigma_{i-1},\sigma_{i+1}]=c_i$. Therefore the internal $\Phi(A_i)$
are the whole edge space. Taking the intersection yields
\[
 C=\gen{c_2,\ldots,c_{d-2}},\qquad |C|=2^{d-3}
       \quad(d\ge3).                                     \tag{5}
\]
At $d=2$, separately, $A_1=\gen{c_1}$ has trivial Frattini subgroup,
so $C=1$.

For every $d\ge4$, $Q_d$ is class two, so $\gamma_3(Q_d)=1$, whereas
(5) makes $C\ne1$. Its order is $2^{2d-1}$, but the quotient by $C$
has order $2^{d+2}$. Thus the source's Theorem 4.2(c), the asserted order
and prescribed-group identification in 4.2(d), and the equality in 4.2(e)
fail on this family. The quotient $Q_d/C$ still has IP: its coordinate
model simply removes the internal edge coordinates, and its parabolics
remain coordinate subgroups. Its interior adjacent product orders become
two. We therefore do not claim to refute every assertion in part (d),
or the separate common-quotient theorem proved above.

\section{Attribution and scope}
The exact edition is arXiv:2605.20691v1, dated 21 May 2026 in the PDF.
Both principal theorems are recovered for their full finite, nondegenerate
range. The explicit auxiliary correction is kept separate from those
positive conclusions. The original common-quotient and covering results
retain the attribution to Hou and Schulte, and the path polytope retains
the source's attribution to Conder. No priority claim is made for the
replacement forest proof or the auxiliary correction.

The full umbrella proof, the marked-group identification in (1)--(2), and
the complete computation (3)--(5) are required proof costs. Existing Lean
components for graded support and kernel lifting do not yet constitute an
end-to-end Lean proof of the forest theorem or this complete application.
No numerical computation or saved group census is a premise here.

\begin{thebibliography}{9}
\bibitem{UA}[UA] \emph{Supported kernels, translation lattices and finite marked quotients},
companion umbrella in this volume.
\bibitem{HS26}[HS26] D.-D. Hou and E. Schulte,
\emph{String C-groups of 2-power order project onto a common string C-group},
arXiv:2605.20691v1, PDF dated 21 May 2026.
\bibitem{Con13}[Con13] M. Conder, \emph{The smallest regular polytopes of
given rank}, \emph{Advances in Mathematics} 236 (2013), 92--110.
\bibitem{MS02}[MS02] P. McMullen and E. Schulte,
\emph{Abstract Regular Polytopes}, Cambridge University Press, 2002.
\end{thebibliography}
\end{document}
